\section{Extensions and Addons}

MIFX allows additional information to be attached to objects through extensions and addons.

This mechanism enables organizations and software vendors to extend the format without modifying the core specification.

Systems that do not recognize an extension or addon must ignore it and continue processing the MIFX package.

\subsection{Extensions}

Extensions provide additional structured metadata stored as separate JSON resources associated with an object.

Extensions are typically stored in an \texttt{extensions} directory adjacent to the object domain.

Example structure:

\begin{verbatim}
setups/
  set-1/
    operations/
      extensions/
        op-1.json
        op-2.json
\end{verbatim}

The meaning of extension fields is defined by the producing system.

Implementations that do not understand an extension must ignore it.

\subsection{Addons}

Addons provide larger or more structured external attachments stored as separate JSON resources.

Addons are placed inside an \texttt{addons} directory associated with a specific object or package domain.

Example structure:

\begin{verbatim}
setups/
  set-1/
    addons/
      setup-addon-1.json

    operations/
      addons/
        op-1.json
        op-2.json
\end{verbatim}

Addons allow complex extensions to be attached without modifying the base object structure.

\subsection{Addon Referencing}

Addons are associated with the object directory in which they reside.

For example:

\begin{verbatim}
setups/set-1/addons/setup-addon-1.json
\end{verbatim}

indicates an addon attached to setup \texttt{set-1}.

Similarly:

\begin{verbatim}
setups/set-1/operations/addons/op-2.json
\end{verbatim}

represents an addon attached to operation \texttt{op-2}.

\subsection{Extension Safety}

Extensions and addons must not alter the interpretation of the core MIFX objects.

Systems must remain able to process the MIFX package even when unknown extensions are present.

This ensures that MIFX remains forward-compatible and vendor-neutral.
